\section{Why pair conductors do not determine quotient capacity}\label{sec:profiles}
\subsection{The parity identity behind the graded calculation}
For the totally ramified complete-order sector, write the successive graded centralizer dimensions as
\[
 n=b_0>b_1>\cdots>b_r=1,
 \qquad \Delta_j=b_{j-1}-b_j.
\]
Assume the actual defining sequence and its order compatibility give descending positive levels $\nu_j$ and the grade-count formulas
\[
 I=\sum_{j=1}^r\Delta_j\left\lceil\frac{\nu_j}{2}\right\rceil-(n-1),
 \qquad
 D=\sum_{j=1}^r\Delta_j\,\mathbf1_{2\mid\nu_j}.
 \tag{24.1}\label{eq:profile}
\]
Here $I=\log_q[I^+:J^1]$, while $D=\dim_{k_F}(J^1/H^1)$. The different floor and ceiling thresholds are essential. In this torsion-one, full-degree setting, the same formal-degree calculation gives
\[
 a(\pi^\vee\times\pi)=n^2-1+(n-1)+2I+D.
 \tag{24.2}\label{eq:pair-profile}
\]
As in the rank-nine calculation, the formula is conditional on actual orders. A list of integers is not a certificate of their arithmetic existence.

\begin{proposition}[Parity cancellation]\label{prop:parity}
Under \eqref{eq:profile},
\[
 (n-1)+2I+D=\sum_{j=1}^r\Delta_j\nu_j.
\]
Consequently
\[
 a(\pi^\vee\times\pi)=n^2-1+\sum_{j=1}^r\Delta_j\nu_j.
 \tag{24.3}\label{eq:pair-weight}
\]
The weighted conductor expression does not retain $D$ as a separate invariant.
\end{proposition}
\begin{proof}
For each integer $\nu$,
\[
 2\left\lceil\frac{\nu}{2}\right\rceil+
 \mathbf1_{2\mid\nu}=\nu+1.
\]
Insert this identity into $2I+D$. Since $\sum_j\Delta_j=n-1$, the constant terms reduce to $-(n-1)$. Adding $n-1$ proves the result.
\end{proof}

This cancellation explains a concrete limitation of the earlier numerical tests. Changing the parity distribution can change the Heisenberg quotient while preserving the weighted sum in \eqref{eq:pair-weight}. More conductor precision does not recover information that the formula has already summed away.

\subsection{The full dyadic profile}
In the inherited full dyadic construction, let $n=2^d$ and let $h_1<\cdots<h_d$ be positive odd integers satisfying its separation requirements. The consecutive profile has
\[
 b_j=2^{d-j},\qquad
 \nu_j=2^{d-j}\,3h_{d-j+1}
 \quad(1\le j\le d).
 \tag{24.4}\label{eq:dyadic-profile}
\]
All levels except the final one are even. Therefore
\[
 D=2^d-2.
 \tag{24.5}\label{eq:dyadic-D}
\]
A source consisting of $d$ nondegenerate binary symplectic planes has dimension $2d$. The original maximal-type quotient is too small for $d=1,2$, has exactly the needed dimension for $d=3$, and has strict surplus for $d\ge4$.

\begin{center}\begin{tabular}{rrrrl}\toprule
$d$ & Arithmetic rank & Source dimension & Host $D$ & Original quotient\\\midrule
1 & 2 & 2 & 0 & Insufficient\\
2 & 4 & 4 & 2 & Insufficient\\
3 & 8 & 6 & 6 & Equal capacity\\
4 & 16 & 8 & 14 & Strict surplus\\
5 & 32 & 10 & 30 & Strict surplus\\\bottomrule
\end{tabular}\end{center}

The rank-two and rank-four enlarged arithmetic carriers in Part I are not counterexamples to the first two rows. Those repairs change the assembly carrier while recovering the same inducing representation. The table concerns the original specified maximal-type quotient. At rank eight, dimension ceases to be the obstruction, but an independently fixed quadratic marking can still obstruct the attachment.

\subsection{Two rank-eight profiles with the same conductor data}
The retained dyadic capacity note contains a sharper example. In the full rank-eight profile, use
\[
 (b_0,b_1,b_2,b_3)=(8,4,2,1),
 \qquad (\nu_1,\nu_2,\nu_3)=(12h_3,6h_2,3h_1).
\]
A second constructed profile skips the intermediate quartic field and has
\[
 (b_0,b_1,b_2)=(8,4,1),
 \qquad (\nu_1,\nu_2)=(12h_3,c),\qquad c=h_1+4h_2.
 \tag{24.6}\label{eq:skip}
\]
The integer $c$ is odd and $2c<12h_3$. Its last jump therefore contributes nothing to $D$. The full profile has $D=6$, while the skipped profile has $D=4$.

Both weighted sums in \eqref{eq:pair-weight} equal
\[
 48h_3+12h_2+3h_1.
 \tag{24.7}\label{eq:sameweight}
\]
Indeed, the full profile has jump sizes $4,2,1$, and the skipped profile has jump sizes $4,3$. The identity $3(h_1+4h_2)=3h_1+12h_2$ supplies the equality.

For $(h_1,h_2,h_3)=(1,3,5)$, the distinction becomes completely explicit.
\begin{center}\begin{tabular}{lrr}\toprule
Quantity & Full profile & Skipped profile\\\midrule
Levels & $(60,18,3)$ & $(60,13)$\\
Principal period & 8 & 8\\
Standard Swan conductor & 60 & 60\\
Standard ordinary conductor & 68 & 68\\
Torsion number & 1 & 1\\
$I$ & 133 & 134\\
$D$ & 6 & 4\\
Heisenberg dimension over $q=2$ & 8 & 4\\
$\delta$ & 279 & 279\\
Ordinary self-pair conductor & 342 & 342\\\bottomrule
\end{tabular}\end{center}
The compensating change is $2\cdot133+6=2\cdot134+4$. The three necessary invariants---standard conductor, torsion number, and self-pair conductor---cannot force the full six-dimensional quotient.

\subsection{The skipped profile is an arithmetic claim, not an integer exercise}
The inherited construction of \eqref{eq:skip} works over $\F_2((t))$. Put $q_0=3h_3$ and take $\alpha^2+\alpha=t^{-q_0}$. A uniformizer $z$ of the quadratic field satisfies
\[
 z=t^{(q_0+1)/2}\alpha,
 \qquad z^2+t^{(q_0+1)/2}z+t=0.
\]
Over that field, the polynomial $X^4+zX+z$ is separable and Eisenstein. If $e$ is a root, let $C$ be multiplication by $e^{-c}$ in the quartic centralizer. The source note chooses a filtered corestriction lift using $s(P_0)=t/z^2$ and takes a correction proportional to $(z^2/t)CP_0$. The actual final parameter is the literal corrected matrix, not an assumed compositum of the approximation fields.

The conclusion that this correction supplies the claimed defining sequence uses the retained refinement argument and its strict inequalities. This publication block has audited the conductor and capacity consequence of that argument. It has not freshly rerun a skipped-profile matrix constructor or recertified every step of its arithmetic field proof. Those statuses are not replaced by the exact integer tests.

Nor is this example a counterexample to H5. Full local pair factors are substantially more information than the three numbers above. The example shows that those three numbers alone do not identify the needed native quotient, and that a source-defined candidate cannot be declared to contain a prescribed type merely because it passes them.

\source{[D167] supplies the actual profile formulas and the skipped-degree construction. Proposition~\ref{prop:parity} and the numerical comparison are fully derived here. The exact suite checks the example and 120 positive odd triples, as well as general parity identities. The underlying arithmetic construction retains its inherited ordinary-proof status.}
